% Vorlage für eine Bachelor- oder Masterthesis
%% Dipl.-Ing. Stefan Averkamp
%% Version 2.5
%% 18.03.2021
%% ----------------------------------------------------------------------------

\documentclass[
cmyk,
pdftex,
a4paper,           % DinA4 Papier
%12pt,             % Schriftgröße
oneside,           % oneside, twoside
openright,         % openleft, openright, openany
numbers=noenddot,  % keinen Punkt hnter Kapitel- oder Sectionnummern
chapterprefix=off, % Kapitel wird vor die Kapitelnummer geschrieben
]{scrbook}
\sloppy
%% Preambel einbinden -----------------------------------------------------------------------------
\UseRawInputEncoding
\setcounter{secnumdepth}{3}         % 3 levels in section numbering

%% Useful packages -------------------------------------------------------------------------------
\usepackage[english]{babel}         % English language
\usepackage{scrhack}
\usepackage{placeins}
\usepackage{subcaption}
\usepackage{lmodern}
\usepackage{cmbright}
\usepackage{enumitem}               % Package for customized lists
\usepackage[OT1]{fontenc}
\usepackage{blindtext}              % Use \blindtext to insert dummy text
\usepackage{mdwlist}                % More compact lists with itemize and co
\usepackage{graphicx}               % Graphics
\usepackage{epstopdf}               % Automatically convert eps images to pdf
\usepackage{ltxtable}               % Combines longtable and tabularx
\usepackage{multirow}               % Multi-row cells in tables
\usepackage[dvipsnames,table,xcdraw]{xcolor}  % Alternating colors in tables
\usepackage{rotating}               % Rotate text and other elements
\usepackage{multicol}               % Multicolumn formatting
\usepackage{amsmath}                % Mathematical formulas
\usepackage{amsthm}                 % Theorem style
\usepackage{mathtools}
\usepackage{amssymb}                % Mathematical symbols
\usepackage{wrapfig}                % Wrapping figures
\usepackage{float}                  % For [h!] or [H] positioning
\usepackage{booktabs}               % For tables from Excel to LaTeX
\usepackage{eurosym}                % Use Euro symbol
\usepackage{url}                    % URLs
\usepackage{etoolbox}
\appto\UrlBreaks{\do\a\do\b\do\c\do\d\do\e\do\f\do\g\do\h\do\i\do\j
    \do\k\do\l\do\m\do\n\do\o\do\p\do\q\do\r\do\s\do\t\do\u\do\v\do\w
    \do\x\do\y\do\z}
\usepackage{adjustbox}
\usepackage{fancybox}               % Boxes and frames
\usepackage{caption}                % Extended settings for captions
\usepackage{array}                  % Extended options for arrays and tables
\usepackage{icomma}                 % Reduces space after comma in decimal numbers
\usepackage{setspace}               % Define line spacing

%% PDF options (links and more) -----------------------------------------------------------------
\usepackage{pdfpages}
\usepackage[utf8]{inputenc}         % UTF-8 encoding
\usepackage[english=british]{csquotes}
\usepackage[pdftex, raiselinks, pdfpagelabels, hypertexnames=false, pdfstartview={Fit}]{hyperref}
\definecolor{lightblue}{rgb}{0.9,0.9,1}

%% PDF metadata
\hypersetup{
    pdfauthor = {Nico Bierbaum},
    pdftitle =  {Scientific Work},
    pdfsubject = {Density and Shrinkage Evaluation of 316L Printed Parts in the FDM Process},
    pdfkeywords = {3D printing, filament, metal, density, tensile test, shrinkage},
    pdfcreator = {LaTeX with hyperref package}
}

%% Header and footer -----------------------------------------------------------------------------
\usepackage[headsepline]{scrlayer-scrpage}
\pagestyle{scrheadings}
\clearpairofpagestyles
\ihead{\headmark} % Redefine
\ohead{\pagemark}

%% Page margins -----------------------------------------------------------------------------------
\usepackage[a4paper,inner=4cm,outer=2.5cm,top=2cm,bottom=2cm,includeheadfoot]{geometry}

%% Bibliography -----------------------------------------------------------------------------------
\usepackage[style=numeric-comp,sorting=none,backend=biber]{biblatex}

\DeclareFieldFormat{labelalpha}{\textsc{#1}}    % Sets the author in small caps 
\renewcommand*{\mkbibnamefamily}[1]{\textsc{#1}}  % Sets the last name in small caps 
\renewcommand*{\mkbibnamegiven}[1]{\textsc{#1}}  % Sets the first name in small caps 
\renewcommand*{\mkbibnameprefix}[1]{\textsc{#1}}  % Sets the prefix in small caps 
\DeclareFieldFormat{title}{``#1''}      % Sets the title in quotation marks
\DeclareFieldFormat{date}{({#1})}               % Sets the date in parentheses
\makeatletter
\def\blx@maxline{77}
\makeatother
\addbibresource{LiteratureDB.bib} % Load bibliography

%% Glossary and abbreviations ---------------------------------------------------------------------
\usepackage{nomencl}
% English title
\renewcommand{\nomname}{Glossary and List of Abbreviations}
% Dots between abbreviation and explanation
\setlength{\nomlabelwidth}{.20\hsize}
\renewcommand{\nomlabel}[1]{#1 \dotfill}
% Reduce line spacing
%\setlength{\nomitemsep}{-\parsep}
\makenomenclature
\setlength{\headsep}{1.5\baselineskip}

%% ---------------------------------------------------------------------------------------------------
%% Create glossary and abbreviations: ------------------------------------------------------------------
%% 1. Run LATEX
%% 2. Open a command prompt and navigate to the folder with the "thesis.tex" file.
%% 3. Run the command   makeindex thesis.nlo -s nomencl.ist -o thesis.nls
%% 4. Run LATEX again. The glossary and abbreviations should now be present.

%% Source code -------------------------------------------------------------------------------------
\usepackage{listings}
\lstset{ %
    basicstyle=\scriptsize,         % the size of the fonts that are used for the code
    numbers=left,                   % where to put the line-numbers
    numberstyle=\footnotesize,      % the size of the fonts that are used for the line-numbers
    backgroundcolor=\color{white},  % choose the background color. You must add \usepackage{color}
    commentstyle=\color{green},
    stringstyle =\color{orange},
    xleftmargin=1.5em,
    xrightmargin=1em,
    frame=tb,                       % adds a frame around the code
    tabsize=2,                      % sets default tabsize to 2 spaces
    captionpos=b,                   % sets the caption-position to top or bottom
    breaklines=true,                % sets automatic line breaking
    breakatwhitespace=false         % sets if automatic breaks should only happen at whitespace
    literate =%
     {#}{\#}1
}
\renewcommand*{\lstlistingname}{Source Code}
\renewcommand*{\lstlistlistingname}{List of Source Codes}

%% Tikz for (complex) colorful drawings ---------------------------------------------------------
\usepackage{tikz}
\usetikzlibrary{positioning,mindmap,shapes,shapes.multipart,shadows,arrows,patterns,topaths}
\usepackage{pgfplots}
\pgfplotsset{width=7cm,compat=newest}

\definecolor{yellow}{HTML}{FFFCCC}
\definecolor{green}{HTML}{CFFFCC}
\definecolor{blue}{HTML}{CCE9FF}
\definecolor{darkblue}{HTML}{99b7EE}
\definecolor{red}{HTML}{FFA8A8}
\definecolor{orange}{HTML}{FFD28F}

\lstset
{
    breaklines=true,
    tabsize=3,
    showstringspaces=false
}

\lstdefinestyle{Common}
{
    extendedchars=true,
    language={[Visual]Basic},
    frame=single,
    %===========================================================
    framesep=3pt,%expand outward.
    framerule=0.4pt,%expand outward.
    xleftmargin=3.4pt,%make the frame fits in the text area. 
    xrightmargin=3.4pt,%make the frame fits in the text area.
    %=========================================================== 
    rulecolor=\color{Red}
}

\lstdefinestyle{A}
{
    style=Common,
    backgroundcolor=\color{Yellow!10},
    basicstyle=\scriptsize\color{Black}\ttfamily,
    keywordstyle=\color{Orange},
    identifierstyle=\color{Cyan},
    stringstyle=\color{Red},
    commentstyle=\color{Green}
}

\lstdefinestyle{B}
{
    style=Common,
    backgroundcolor=\color{Black},
    basicstyle=\scriptsize\color{White}\ttfamily,
    keywordstyle=\color{Orange},
    identifierstyle=\color{Cyan},
    stringstyle=\color{Red},
    commentstyle=\color{Green}
}
% Default fixed font does not support bold face
\DeclareFixedFont{\ttb}{T1}{txtt}{bx}{n}{12} % for bold
\DeclareFixedFont{\ttm}{T1}{txtt}{m}{n}{12}  % for normal

% Custom colors
\usepackage{color}
\definecolor{deepblue}{rgb}{0,0,0.5}
\definecolor{deepred}{rgb}{0.6,0,0}
\definecolor{deepgreen}{rgb}{0,0.5,0}

% Python style for highlighting
\newcommand\pythonstyle{\lstset{
language=Python,
basicstyle=\ttm,
morekeywords={self},              % Add keywords here
keywordstyle=\ttb\color{deepblue},
emph={MyClass,__init__},          % Custom highlighting
emphstyle=\ttb\color{deepred},    % Custom highlighting style
stringstyle=\color{deepgreen},
frame=tb,                         % Any extra options here
showstringspaces=false
}}


% Python environment
\lstnewenvironment{python}[1][]
{
\pythonstyle
\lstset{#1}
}
{}
\begin{document}

%% Titelseite einbinden ---------------------------------------------------------------------------
\thispagestyle{empty}
\newgeometry{
    top=25mm, 
    bottom=7mm,
    left=46mm,
    right=36mm
}
%Logos------------------------------------------------------------------------------
%Logos have a defined width of 6cm, the box containing the logo is 7cm
\vspace*{-2cm} % may be useful to place logos higher
\setlength{\unitlength}{1cm}
\hfill
\begin{minipage}[t]{4cm} 
    \begin{center}
       \includegraphics[width=40mm]{bilder/Firmenlogo.png}
    \end{center}
\end{minipage}
 %% University Logo ----------------------------------------------------------------------
 \hfill  % This adds space and pushes the next minipage to the right
 \begin{minipage}[t]{7cm}
    \begin{flushright}
       \includegraphics[width=66mm]{bilder/fhm_Logo_CMYK_30mm.pdf}
    \end{flushright}
 \end{minipage}
%-----------------------------------------------------------------------------------

\begin{minipage}[t]{11cm}
\vspace{2cm}                    
\begin{Huge}
\begin{center}\textsf{\textbf{Master's Thesis}}\end{center}
\end{Huge}


\vspace{2cm}                                
\begin{Large}
\hfill 
\begin{center}

\begin{large}
    Nico Bierbaum
\end{large}

\vspace{2cm} %3cm
%\begin{minipage}[t][2.5cm]{14cm}
\begin{flushright}
    \hfill{Development of a Robot-Based Method for Characterizing Single-Use Bag Mixers}\\
\end{flushright}            
%\end{minipage}
\end{center}    
\end{Large}

\begin{large}
\vspace{30mm}  %3cm
%\begin{minipage}[t][1.5cm]{14cm}
\hfill Supervisor: Prof. Dr.-Ing. Stephan Behr

\hfill Marius Kling M.Sc.\\
%\end{minipage}
\end{large}

\vspace{35mm}
\hfill Leverkusen, \today
\end{minipage}
\vspace{20mm}

\begin{minipage}[t]{136.5mm}
\begin{flushright}
\includegraphics[width=82mm]{bilder/MB_CMYK.pdf}
\end{flushright}
\end{minipage}
\restoregeometry
%--------------------------------------------------------------------------------------------------------


%% Sperrvermerk, Eidesstattliche Erklärung, Vorwort -----------------------------------------------
\frontmatter

\pagestyle{scrheadings}
\pagenumbering{Roman}


\input{inhalt/Sperrvermerk.tex}
%Eidesstattliche Erklärung-------------------------------------------------------------------------------
\chapter*{Affidavit}
\addcontentsline{toc}{chapter}{Affidavit}

I hereby declare under oath that I have completed this term paper independently and without unauthorized external assistance. The literature sources used are fully cited in the bibliography.

\vspace{6cm}
Leverkusen, \today

% Preface -----------------------------------------------------------------------------------------
\chapter*{Preface}

This Master's thesis presents the development and implementation of a robot-based methodology for the characterization of Single-Use Bag Mixers. The aim of this research is to advance the precision and efficiency of mixing performance assessments within bioprocess engineering.

The project addresses critical challenges in the industry, including the necessity for precise and reproducible evaluations of mixer performance and the integration of automated systems in laboratory environments. The developed robotic method not only improves the accuracy of these assessments but also significantly enhances the efficiency by reducing the time and manual effort required.

This research has been conducted under the supervision of Prof. Dr.-Ing. Stephan Behr, whose expert guidance and insights were invaluable throughout the course of this project. I am also deeply grateful to Marius Kling, M.Sc., for his technical support and constructive feedback.

I would like to acknowledge the support and collaboration of [Company/Institution Name], whose resources and partnership were instrumental to the success of this project.

The findings and methodologies presented in this thesis contribute to the advancement of bioprocessing technologies, providing a robust foundation for future research and development in the field.

Leverkusen, \today

\addcontentsline{toc}{chapter}{Preface}


%% Inhaltsverzeichnis einbinden -------------------------------------------------------------------
\pdfbookmark[0]{Inhaltsverzeichnis}{tableofcontents}
\tableofcontents

%% Verzeichnisse einbinden  -----------------------------------------------------------------------
\listoffigures              %% Abbildungsverzeichnis
\listoftables               %% Tabellenverzeichnis
\lstlistoflistings          %% Quelltextverzeichnis
\input{inhalt/glossary.tex}  %% Abkuerzungsverzeichnis
%% Eigentlicher Inhalt der Arbeit -----------------------------------------------------------------
%% Einbinden der einzelnen Kapitel. Jedes Kapitel in eine seperate Datei speichern. ---------------
\mainmatter
\chapter{Introduction}

\begin{itemize}
    \item This work is based on VDI VDE 2206
    \item 
\end{itemize}
\chapter{Technical Background}
\section{Development of a cyber-physical system}

A Cyber-Physical System (CPS) integrates computation, networking, and physical processes. Embedded computers and networks monitor and control physical processes, usually with feedback loops where physical processes affect computations and vice versa. CPS is critical in applications like smart grids, autonomous vehicles, medical monitoring, industrial control systems, and robotics. Key characteristics include real-time operation, robustness, and integration of hardware and software components to ensure seamless interaction between the cyber and physical elements.
\autocite{Lee2008}
\section{Methods for the characterization of (bag) mixing systems}
\subsection{Characterization aspects \- power input, suspension, mixing time, shear stress}
\subsection{Evaluation methods \- including image-based evaluation methods}
\section{Robot-based (bag) mixing systems}
\input{inhalt/3_Requirements_analysis.tex}
\chapter{System Design}

\section{List of Symbols}


\begin{table}[htbp]
    \centering
    \begin{tabular}{l l}
        \hline
        \textbf{Symbol/Formula} & \textbf{Description} \\
        \hline
        \((KS)_E\) & Origin of the coordinate system on the inclined plane \\
        \( _{(E)}\vec{p}(t) \) & 3D vector of the total motion in plane \((KS)_E\)\\
        \(_{(E)}\vec{p_B} (t) \) & 3D vector of pure motion in plane \((KS)_E\) \\
        \(_{(E)}\vec{p_0}\) & 3D start vector. User input (gap width)\\
        \(z(t)\) & Motion equation of the robot trajectory \\
        \(\hat{z}\) & User input for the lifting height\\
        \(\omega\) & Angular frequency \\
        \(f\) & User input for the frequency \\
        \(W\) & Physical work \\
        \(F\) & Physical force \\
        \(P\) & Physical power \\
        \hline
    \end{tabular}
    \caption{List of Symbols and Formulas}
\end{table}

\section{Motion Profile of the Paddle Mixer}

\begin{figure}[h]
    \centering
    \includegraphics[width=0.8\textwidth]{bilder/KS_E-Koordinatensystem_Ebene.png}
    \caption[Visualization of the Coordinate System on the Inclined Plane]{Visualization of the Coordinate System on the Inclined Plane (own illustration)}\label{fig:VisualizationCoordinateSystemInclinedPlane}
\end{figure}

\noindent The following describes the three parameters for a reproducible description of the paddle movement in all three directions of motion. Here, \( _{(E)}\vec{p}(t) \) represents the total movement in the coordinate system \((KS)_E\). \(_{(E)}\vec{p_B} (t) \) describes the variable position starting from \(_{(E)}\vec{p_0}\).\\

\textbf{General:}

\begin{equation}
    \begin{aligned}
        _{(E)}\vec{p}(t) &= _{(E)}\vec{p_B}(t) + _{(E)}\vec{p_0}\\
        \text{with}\\
        _{(E)}\vec{p}_B(t) &= _{(E)}\left(x_B(t), y_B(t), z_B(t)\right)\\
        _{(E)}\vec{p_0} &= _{(E)}\left(x_0, y_0, z_0\right)
    \end{aligned}
    \label{eq:generalMovement}
\end{equation}

\autoref{eq:generalMovement} is independent of the initiated type and direction of movement. Any possible types of movement can be represented in \(\vec{p}_B(t)\).\\

The calculation based on a sinusoidal movement along the Z-axis follows:

\begin{equation}
    \begin{aligned}
        _{(E)}\vec{p}_{0} &= (x_{0}, y_{0}, z_{0})\\
        _{(E)}\vec{p}(t) &= (x, y, z(t)) + \vec{p}_{0}  
    \end{aligned} 
\end{equation}

\( z(t) \) is defined by:

\begin{equation}
    z(t) = -\hat{z}(\cos(\omega t)-1)  \quad \text{with} \quad \hat{z} = \frac{h}{2}
\end{equation}

and \( \omega = 2 \pi f = \frac{2 \pi}{T} \).

For example, at \( t = 0 \):

\begin{equation}
    z(0) = -\hat{z}(1-1) = 0
\end{equation}

Example at \( t = \frac{1}{4}T \):

\begin{equation}
    \begin{aligned}
        z\left(\frac{1}{4}T\right) &= -\hat{z}\left(\cos\left(\frac{2\pi T}{T \cdot 4}\right) - 1\right) \\
        &= -\hat{z}\left(\cos\left(\frac{\pi}{2}\right) - 1\right) \\
        &= -\hat{z}(0 - 1) = \hat{z} = \frac{1}{2} h
    \end{aligned}
\end{equation}

And at \( t = \frac{1}{2}T \):

\begin{equation}
    \begin{aligned}
        z\left(\frac{1}{2}T\right) &= -\hat{z}\left(\cos\left(\frac{2\pi T}{T \cdot 2}\right) - 1\right) \\
        &= -\hat{z}\left(\cos(\pi) - 1\right) \\
        &= -\hat{z}(-1 - 1) = 2\hat{z} = h
    \end{aligned}
\end{equation}

where \( _{(E)}\vec{p(t)} \) represents the movement of the paddle center and \( _{(E)}\vec{p}_{0} \) represents the center on the single-use bag.

\begin{figure}[h]
    \centering
    \includegraphics[width=0.8\textwidth]{bilder/cosine-plot.png}
    \caption[Cosine Function]{Classical cosine function (\( cos(t) \)) (own illustration)}\label{fig:CosineFunction}
\end{figure}
\FloatBarrier

\begin{figure}[h]
    \centering
    \includegraphics[width=0.8\textwidth]{bilder/cosine-plot_real.png}
    \caption[Representation of Paddle Movement Using the Cosine Function]{Representation of the paddle movement in the Z-direction using the cosine function \( z(t) = -\hat{z}(\cos(\omega t)-1) \) (own illustration)}\label{fig:PaddleMovementCosine}
\end{figure}
\FloatBarrier

Example sinusoidal movement in \( _{(E)}\) Z-direction:
\begin{equation}
    \begin{aligned}
        x &= x_B + x_0\\
        y &= y_B + y_0\\
        z(t) &= \hat{z}\left( 1 - \cos(\omega t) \right) +z_0\\
        \text{with}\\
        &\hat{z} = 2h\\
        &\omega = 2\pi f
    \end{aligned}
\end{equation}

\section{Definition of Power Input for the Mixing Process}

Generally, it holds:

\begin{equation}
    W = \int_{t_1}^{t_2} \vec{F}(\vec{s(t)})\cdot d \vec{s}
\end{equation}

\noindent If the force direction matches the motion direction and both vectors point in the same direction, it holds:

\begin{equation}
    W = F \cdot s \text{   and   } P = \frac{W}{dt}
\end{equation}

\noindent In our case, the force \(F_z\) points in the direction of motion \(z\), but it is not constant. Therefore, it holds:

\begin{equation}
    \begin{aligned}
        W = \int_{t_1}^{t_2} F_z (z(t))\cdot dz\\
        W = \int_{t_1}^{t_2} F_z (t)\cdot dz
    \end{aligned}
\end{equation}

The relationship between force and distance is

\chapter{Component Design}

In this chapter the system design out of \autoref{sec:SystemDesign} gets applied. 

\section{Software}

Software is a fundamental part of the design of this thesis. To combine all components (e.g. Camera, force sensor, report, HMI) a software is needed. 

\subsection{Parameterizable Robot Program for Executing the Movement}



\subsection{Development of an Automated Evaluation Script}
\subsubsection{For Evaluating the Power Input}

In \autoref{sec:Motion Profile of the Paddle Mixer}, the theoretical background about the calculation of the input power is described. In this subsubsection the equations are translated to python code.

\begin{python}
    class MyClass(Yourclass):
        def __init__(self, my, yours):
            bla = '5 1 2 3 4'
            print bla
\end{python}
\subsubsection{For Evaluating the Mixing Time or Suspension (e.g., image-based)}
\subsubsection{With Easy Usability (for Laboratory Users) e.g., with GUI/Tkinter}
\subsubsection{Generation of CSV Data and Reports}
\section{Hardware}
\subsection{Construction of an Inclined, Angle-Adjustable Support Surface for the Bag}
\subsection{Construction of Sensor Systems to Determine Mixing Time or Suspension (e.g., Camera + Lighting under the Bag)}

\subsubsection{Integration of Sensor and Camera Systems}

\paragraph{Payload and Center of Gravity Calculation}

The integration of the SCHUNK FT-AXIA 80 sensor and the Robotiq wrist camera requires accurate calculations for the payload and center of gravity to ensure precise and stable operation of the robotic system. The following sections detail these calculations.

\subparagraph{Payload Calculation}

The payload is calculated as the total mass of the sensor, camera, and custom tool:

\begin{equation}
\text{Payload} = \text{Mass of sensor and adapter} + \text{Mass of camera} + \text{Mass of custom tool}
\end{equation}

Given:
\begin{align*}
\text{Mass of sensor and adapter} &= 0.558 \, \text{kg} \\
\text{Mass of camera} &= 0.23 \, \text{kg} \\
\text{Mass of custom tool} &= 0.16243 \, \text{kg}
\end{align*}

\begin{equation}
\text{Payload} = 0.558 + 0.23 + 0.16243 = 0.95043 \, \text{kg}
\end{equation}

\subparagraph{Center of Gravity Calculation}

To find the center of gravity (CG) in the Z direction, the combined CG of all components (sensor, camera, and custom tool) is calculated.

\textbf{Sensor}:
\begin{align*}
\text{Mass} &= 0.558 \, \text{kg} \\
\text{CG in Z direction} &= 22.8 \, \text{mm}
\end{align*}

\textbf{Camera}:
\begin{align*}
\text{Mass} &= 0.23 \, \text{kg} \\
\text{Total offset in Z direction} &= 23.5 \, \text{mm} + 9 \, \text{mm} + 14.5 \, \text{mm} = 47 \, \text{mm}
\end{align*}

\begin{equation}
\text{Combined CG}_{\text{sensor+camera}} = \frac{0.558 \times 22.8 + 0.23 \times 47}{0.558 + 0.23}
\end{equation}

\begin{equation}
\text{Combined CG}_{\text{sensor+camera}} = \frac{12.7224 + 10.81}{0.788} = \frac{23.5324}{0.788} = 29.85 \, \text{mm}
\end{equation}

\textbf{Custom Tool}:
\begin{align*}
\text{Mass} &= 0.16243 \, \text{kg} \\
\text{Total offset in Z direction} &= 30 \, \text{mm} + 18.596 \, \text{mm} = 48.596 \, \text{mm}
\end{align*}

\begin{equation}
\text{Combined CG}_{\text{system}} = \frac{(0.788 \times 29.85) + (0.16243 \times 48.596)}{0.95043}
\end{equation}

\begin{equation}
\text{Combined CG}_{\text{system}} = \frac{23.5282 + 7.8981}{0.95043} = \frac{31.4263}{0.95043} = 33.07 \, \text{mm}
\end{equation}

\subparagraph{Summary}
\begin{itemize}
    \item \textbf{Payload}: 0.95043 kg
    \item \textbf{Center of Gravity (Z)}: 33.07 mm
\end{itemize}

These values incorporate the mass and CG offsets of both the camera and the custom tool, ensuring accurate configuration for the robot setup.
\chapter{Implementation (Optional)}
\section{Construction and Integration of the Complete System}

\chapter{Validation/Evaluation}
\section{Execution and Documentation of Tests to Validate System Functions According to the Requirements from Chapter 3}
\section{Validation of System Functions}

\input{inhalt/8_Conclusion_Outlook.tex}

%% ...
%%\chapter{Diskursiver Teil}
\label{Zusammenfassung}



%% Beginn vom Anhang ------------------------------------------------------------------------------
\appendix

%% Literaturverzeichnis einbinden. Die Quellen sind in der Datei "LiteraturDB.bib" gespeichert. ---
%% Das Literaturverzeichnis muss vorab über über Tools-->Befehle-->Biber erzeugt werden.--
\begin{flushleft}  % Schaltet den Blocksatz für das Literaturverzeichnis ab.
     \printbibliography
\end{flushleft}
%% Anhang einbinden -------------------------------------------------------------------------------
\input{inhalt/appendix.tex}

\end{document}
